% Pillar-7 synthesis: what the Pillar-2 cross-examination predicted, what held here
\section{Synthesis: What the Pillar-2 Cross-Examination Predicted, and What
Held}
\label{sec:p7-synthesis}
The companion Pillar-2 paper closed with an external frontier-model
cross-examination that produced, as \emph{reasoning contributions},
several falsifiable predictions about ZVF. This paper is, in part, the
empirical answer to that cross-examination, and it is worth scoring the
predictions explicitly.

\paragraph{Predicted and confirmed.} (1) \emph{The binomial form.} The
cross-examination proposed $\mathrm{ZVF}(p, G) = p^G + (1-p)^G$ as an
inelasticity result; our T1 adopts it, and the endpoint/interior profile of
the binned reward tensors plus the $h_G(\bar p)$ column of
Table~\ref{tab:p7-gsweep} bear it out. (2) \emph{Mastery--incapacity
aliasing as the structural defect.} Predicted from the symmetry
$h_G(p) = h_G(1-p)$; confirmed at population scale by the 368-run U-shape
(Table~\ref{tab:p7-ushape}), which is the aliasing rendered as data.
(3) \emph{The jitter fragility.} The cross-examination's sharpest sub-test
was that sub-reward jitter would zero ZVF without changing learning; the
$0.158 \to 0.000$ collapse with PCD invariant
(Section~\ref{sec:p7-jitter}) confirms it exactly. (4) \emph{Magnitude-aware
replacement.} The proposed pairwise-contrast density is implemented here as
PCD, with its expectation $\frac{G-1}{G}\E[p(1-p)]$ derived and its parabola
measured.

\paragraph{Predicted and still open.} The cross-examination set a bar of
$\rho \geq 0.45$ for any replacement diagnostic as an \emph{early-window
predictor of held-out outcome}. We have not cleared that bar, and we did
not try to clear it with proxies: the honest test needs per-group reward
tensors on every anchor run, which only the GSM8K runs carry. What we
\emph{can} report is the sign-resolution half of the claim: on the 80-run
summary corpus, the sign-carrying statistic (mean reward) reaches
$\rho = 0.95$ against outcome where unsigned ZVF reaches $0.56$, consistent
with the cross-examination's diagnosis that the missing ingredient is sign,
not more of the same telemetry.

\paragraph{Beyond the predictions.} Two results here were not anticipated
by the cross-examination. First, the interventional turn: T2 is not only an
explanation but a control law, and E3 shows a live controller acting on it
matches the best fixed recipe on held-out gain. Second, the
accuracy-versus-compute framing of dynamic sampling: the cross-examination
treated ZVF $= 0$ as a target, whereas the E3 audit prices it ($+45\%$
rollouts) and finds the purchase optional on a well-conditioned task. The
program-level lesson for the pillar family is the same one Papers~1--4
reached from four different directions: telemetry must be
magnitude-aware, stack-attributed, and priced---or it will be gamed by the
pipeline, the backend, or the budget.
